%\pagebreak
\section{\kcode{fuse} Construct}
\label{sec:fuse}
\index{constructs!fuse@\kcode{fuse}}
\index{fuse construct@\kcode{fuse} construct}
\index{looprange!full@\kcode{looprange}}
\index{looprange clause@\kcode{looprange} clause}
\index{clauses!looprange@\kcode{looprange}}

The \kcode{fuse} construct applies loop fusion to a loop nest sequence, as
specified by a \kcode{looprange} clause. 
The first argument of the \kcode{looprange} clause specifies the first 
loop nest to merge in a sequence of canonical loop nests that follows,
and the second argument specifies the number of loops to merge.
Without the \kcode{looprange} clause, the merge applies to the entire sequence.

The effect of fusion is to merge 
separate iteration spaces of a sequence of loop nests
into a single iteration space of a single canonical loop
with an iteration count of the largest loop and the loop blocks
combined in program order. Other loop transforming constructs
may be applied to the results of a \kcode{fuse} construct.

Several examples below illustrate the application of the \kcode{fuse} construct
and show equivalent fused loops that illustrate the functionality of the constructs.
The equivalent routines may undergo further compiler optimization.


In the following example sine and cosine values are created for
a list of identical values in separate loops.
By fusing the loops, each iteration of the fused loop calculates
a sine-cosine pair having the same argument value.
Under this condition compilers can often calculate the pair
in significantly less time.

\cexample[6.0]{fuse}{1}
\ffreeexample[6.0]{fuse}{1}


The following example illustrates the fusion of two loops with different lengths.
Here the fusion merges the two
separate iteration spaces (\ucode{i} and \ucode{j}) into a single 
logical iteration space (\ucode{ij})
with a size of the largest loop, and their loop blocks
combined with conditionals to avoid 
execution of out-of-range iterations for the shorter loop.
The fusion is applied to minimize loop overhead, and possibly use
the caches more efficiently when the function \ucode{f()} and \ucode{g()} 
are computed in the same iteration.
Fusion of different-sized loops works for C++11 range-based for loops.

\cexample[6.0]{fuse}{2}
\ffreeexample[6.0]{fuse}{2}

In the following example the \kcode{looprange} clause is used to select two loops
(\ucode{2} in the second argument)
from the loop nest sequence for fusion into a canonical loop. 
Loop fusion begins with the second loop (\ucode{2} in the first argument)
of the loop nest sequence. All the other loops are unaltered.

\cexample[6.0]{fuse}{3}
\ffreeexample[6.0]{fuse}{3}


